% Unite: papier 34, section 6; traduction francaise depuis la source allemande auditee.

\providecommand{\mG}{\mathfrak{G}}
\providecommand{\mE}{\mathfrak{E}}
\providecommand{\mH}{\mathfrak{H}}
\providecommand{\mR}{\mathfrak{R}}
\providecommand{\mo}{\mathfrak{o}}

\subsection*{\S{} 6. Modules relativement à un corps. Systèmes hypercomplexes}

Un exemple presque trivial de groupes avec opérateurs complètement réductibles est fourni par les modules ayant une base finie relativement à un corps (non nécessairement commutatif).

Soit $\mG$ un module à droite sur $K$, et soit l'élément unité $e$ de $K$ en même temps opérateur unité: $ae=a$ pour $a$ dans $\mG$. Tout sous-module $aK$ dérivé d'un $a$ est simple, c'est-à-dire égal au module dérivé de l'un quelconque de ses éléments $\ne0$. Si donc $\mH$ est le sous-module dérivé de certains éléments, et si $a$ est un élément qui n'appartient pas à $\mH$, alors $\mH\cap aK=\mE$, donc $\mH+aK$ est directe. Ainsi, en partant d'un élément de base quelconque $a_1$, on peut toujours adjoindre de nouveaux éléments de base $a_2,a_3,\ldots$, et l'on obtient $\mG$ comme somme directe
\[
        \mG=a_1K+a_2K+\cdots+a_nK.
\]
\emph{Donc $\mG$ est complètement réductible.}

Le nombre $n$, longueur de la suite de composition, s'appelle le \emph{rang} de $\mG$ relativement à $K$. En raison de la représentation univoque des éléments de $\mG$ (et parce que $e$ a été supposé opérateur unité), la base $a_1,\ldots,a_n$ est linéairement indépendante.

Tout sous-module $\mH$ est un sommant direct:
\[
        \mG=\mH+\mR=(h_1K+\cdots+h_sK)+(r_1K+\cdots+r_{n-s}K),
\]
ou encore: \emph{on peut compléter toute base linéairement indépendante de $\mH$ en une base linéairement indépendante de $\mG$.} Les $r_i$ peuvent même être choisis parmi les éléments de base initiaux $a_i$.

Soient $(a_1,\ldots,a_n)$ et $(c_1,\ldots,c_n)$ des bases linéairement indépendantes. Alors
\[
        c_k=\sum_i a_i\pi_{ik},
\]
ou, écrit en matrices:
\[
        (c_1,\ldots,c_n)=(a_1,\ldots,a_n)P,
        \qquad
        P=\begin{pmatrix}
        \pi_{11}&\cdots&\pi_{1n}\\
        \vdots&&\vdots\\
        \pi_{n1}&\cdots&\pi_{nn}
        \end{pmatrix}.
\]
Réciproquement,
\[
        (a_1,\ldots,a_n)=(c_1,\ldots,c_n)Q.
\]
Il s'ensuit
\[
        (a_1,\ldots,a_n)=(a_1,\ldots,a_n)PQ,
        \qquad
        (c_1,\ldots,c_n)=(c_1,\ldots,c_n)QP,
\]
donc, puisque les $a_i$ et les $c_i$ sont linéairement indépendants,
\[
        PQ=QP=E.
\]

Les $K$-modules particuliers qui sont en même temps des anneaux sont les ``systèmes hypercomplexes''.

Un anneau $\mo$, qui est en même temps un module à droite relativement à un corps commutatif $K$, s'appelle \emph{système hypercomplexe relativement à $K$} (également désigné dans la littérature comme ``algèbre sur $K$''), si:
\begin{enumerate}
\item le rang est fini (soit $u_1,\ldots,u_n$ une base linéairement indépendante);
\item $ab\cdot x=a\cdot bx=ax\cdot b$. Autrement dit, les actions de $K$ et de $\mo$ sur le module commutent;
\item l'élément unité $e$ de $K$ est en même temps opérateur unité: $ae=a$ pour $a$ dans $\mo$.
\end{enumerate}
Comme
\[
        \Bigl(\sum_i u_i\alpha_i\Bigr)\Bigl(\sum_k u_k\beta_k\Bigr)
        =\sum_{i,k}(u_i u_k)(\alpha_i\beta_k)
\]
le système est déterminé univoquement lorsque, outre $K$, est donnée la \emph{table de multiplication}, qui indique comment chaque produit $u_i u_k$ s'exprime au moyen des $u_j$:
\[
        u_i u_k=\sum_j u_j\gamma^j_{ik}.
\]
Les $\gamma^j_{ik}$ doivent satisfaire aux relations connues qui résultent de la loi associative.

Si $\mo$, comme nous le supposerons souvent, possède un élément unité $e$ ($ea=ae=a$ pour tout $a$), alors on peut identifier les éléments $x$ de $K$ avec $ex$, donc considérer $K$ comme un sous-corps de $\mo$. Le système hypercomplexe peut alors aussi être décrit comme \emph{un anneau de rang fini relativement à un corps situé dans le centre.}

Un exemple de système hypercomplexe est l'\emph{anneau de groupe} d'un groupe fini; il apparaît lorsqu'on prend pour éléments de base $u_i$ les éléments d'un groupe fini, et pour multiplication la multiplication du groupe. Le corps $K$ est ici arbitraire.

\footnotetext{La condition 2. exprime que la multiplication par un élément de $K$ signifie un homomorphisme pour $\mo$, lorsque $\mo$ est considéré aussi bien comme $\mo$-module à gauche que comme $\mo$-module à droite. En vertu de 3., ces homomorphismes sont même des isomorphismes.}
